\makeatletter
\def\input@path{{./}{../../../preamble/}{book2026-2027/preamble/}}
\makeatother
\input{preamble.tex}
\input{graphs/highlights.tex}
\GTenHighlightsfalse % Uncomment to disable instructional highlights.
\input{graphs/specific_highlights.tex}
\CoordinatedTermTwoPalette
\usepackage{needspace}
\hypersetup{pdftitle={Grade 10 Term 2 Discrete Probability}}
\setlength{\emergencystretch}{2em}
\title{Grade 10, Term 2, Lesson 1, Criterion C4\\Binomial Counts in Fixed Financial Trials}
\author{}
\date{}
\begin{document}
\maketitle
\setcounter{tocdepth}{1}
\setcounter{secnumdepth}{2}
\phantomsection\label{toc}
\tableofcontents
\section*{Evaluated criteria}
\begin{description}
\item[\textbf{C4}] Applies the binomial distribution to fixed-trial financial situations by verifying its conditions and calculating probabilities using the binomial probability formula.
\end{description}
\textbf{How to practise.} Attempt each C4 task before following its solution link. Define success, check all four binomial conditions and translate the event before calculating. Keep exact intermediate values and round final probabilities to four decimal places. Use the arrows to return to the contents, problem or problem group.
\clearpage
\section{Introduction}\label{sec:t2-c4-introduction}
\subsection*{Highlight legend}
\phantomsection\label{tab:highlight-legend}
\PracticeHighlightLegend{G10-T2-C4-coordinated}
Names, notation and units identify the roles even without color. Related shades connect related quantities; their definitions distinguish them.

All data in this activity are hypothetical teaching models. Probabilities have no monetary units. Keep exact expressions where useful and round final probabilities to four decimal places.
\subsection{Introductory example 1: a fixed set of offers}\label{ex:intro-one}
A school shop sends four independent offers. Each visitor either purchases (success) or does not, with constant purchase probability $\HLSuccessProbability{p}=\HLSuccessProbability{0.3}$. Let $\GTenHLValue{X}$ count \GTenHLValue{purchases}. The number of trials $\HLObservationCount{n}=\HLObservationCount{4}$ is fixed before results are observed. The support is $\{0,1,2,3,4\}$ purchases. Probability has no monetary units; revenue might be a separate variable $\GTenHLSecond{Y=5X}$ USD.
\textbf{Check.} Does the shop stop after its first sale? \emph{Suggested answer:} No. A binomial count observes all four fixed trials.


\clearpage
\subsection{Introductory example 2: a fixed set of payment reminders}\label{ex:intro-two}
A business sends \HLObservationCount{three reminders}, each concerning a different invoice. Assume independent results and constant probability $\HLSuccessProbability{p}=\HLSuccessProbability{0.5}$ that each invoice is paid by Friday. Let $\GTenHLValue{X}$ count paid invoices, with support $\GTenHLValue{\{0,1,2,3\}}$ invoices. Exactly $\GTenHLValue{k}=\GTenHLValue{2}$ payments can appear in three orders: SSF, SFS and FSS. Here S means paid by Friday and F means not paid by Friday. Each order has probability $\HLSuccessProbability{0.5}\cdot\HLSuccessProbability{0.5}\cdot\HLComplementProbability{0.5}$; adding the three disjoint orders gives $\GTenHLProbability{0.375}$. This is a probability, not a dollar amount.
\textbf{Quick check.} If the firm stops sending reminders as soon as an invoice is paid, is the number of reminders fixed? \emph{Suggested answer:} No; a fixed-trial binomial count no longer describes that stopping rule.

\subsection{Four conditions}
Use binomial only with a fixed number $n$ of binary trials, independent results and constant success probability $p$. Independent means one result does not change the probability of another. A common economic shock may break independence. Changing customer groups may break constant probability.
\textbf{Check.} Are default events automatically independent? \emph{Suggested answer:} No. A shared economic shock can link them.


\subsection{Count the arrangements}
Exactly two successes in three trials can occur as SSF, SFS or FSS. Each sequence has probability $p^2(1-p)$. There are three such arrangements, so add their probabilities. In general their number is the combination $\binom nk$, read as n choose k:
\[\binom nk=\frac{n!}{k!(n-k)!}.\]
The factorial $n!$ means multiply positive integers from 1 to $n$: $4!=4\cdot3\cdot2\cdot1=24$, and $0!=1$. Thus $\binom42=24/(2\cdot2)=6$. Order matters for a sequence, but the combination counts all sequences with the same number of successes.
\textbf{Check.} Why is $p^2(1-p)$ alone insufficient for two successes in three trials? \emph{Suggested answer:} It describes only one arrangement; there are three.

\textbf{Check.} What does $\sum_{j=0}^{2}P(X=j)$ mean? \emph{Suggested answer:} Add $P(X=0)+P(X=1)+P(X=2)$; it is an at-most probability.

\clearpage
\section{Notation Reference and Practice Procedure}\label{sec:notation-and-procedure}
\subsection{Notation reference}\label{tab:notation-reference}
\begingroup
\setlength{\tabcolsep}{4pt}
\begin{longtable}{@{}>{\raggedright\arraybackslash}p{.16\linewidth}>{\raggedright\arraybackslash}p{.18\linewidth}>{\raggedright\arraybackslash}p{.27\linewidth}>{\raggedright\arraybackslash}p{\dimexpr.39\linewidth-6\tabcolsep\relax}@{}}
\toprule
\textbf{Symbol or notation} & \textbf{Concept} & \textbf{Definition} & \textbf{Meaning in the introductory examples} \\
\midrule\endfirsthead
\toprule
\textbf{Symbol or notation} & \textbf{Concept} & \textbf{Definition} & \textbf{Meaning in the introductory examples} \\
\midrule\endhead
$\GTenHLValue{X}$; $\GTenHLValue{k}$ & \GTenHLValue{Success count} & $X$ is the random count; $k$ is a possible value, measured in successes. & \textbullet\ \hyperref[ex:intro-one]{Example 1}: purchases.\par\textbullet\ \hyperref[ex:intro-two]{Example 2}: paid invoices; $\GTenHLValue{k}=\GTenHLValue{2}$ specifies exactly two. \\[4pt]
$\GTenHLValue{\{0,\ldots,n\}}$ & \GTenHLValue{Support} & All possible integer success counts, from zero to the fixed trial count. & \textbullet\ \hyperref[ex:intro-one]{Example 1}: $\GTenHLValue{\{0,1,2,3,4\}}$ purchases.\par\textbullet\ \hyperref[ex:intro-two]{Example 2}: $\GTenHLValue{\{0,1,2,3\}}$ paid invoices. \\[4pt]
$\HLObservationCount{n}$ & \HLObservationCount{Trial count} & Number of trials fixed before observing outcomes. & \textbullet\ \hyperref[ex:intro-one]{Example 1}: $\HLObservationCount{4}$ offers.\par\textbullet\ \hyperref[ex:intro-two]{Example 2}: $\HLObservationCount{3}$ reminders. \\[4pt]
$\HLSuccessProbability{p}$ & \HLSuccessProbability{Success probability} & Constant probability of success on one trial; $0\leq p\leq1$. & \textbullet\ \hyperref[ex:intro-one]{Example 1}: purchase probability $\HLSuccessProbability{0.3}$.\par\textbullet\ \hyperref[ex:intro-two]{Example 2}: payment probability $\HLSuccessProbability{0.5}$. \\[4pt]
$\HLComplementProbability{1-p}$ & \HLComplementProbability{Failure probability} & Probability of the other binary outcome. & \textbullet\ \hyperref[ex:intro-one]{Example 1}: no purchase, $\HLComplementProbability{0.7}$.\par\textbullet\ \hyperref[ex:intro-two]{Example 2}: unpaid by Friday, $\HLComplementProbability{0.5}$. \\[4pt]
$\binom nk$; $n!$ & Arrangements; factorial & $\binom nk=n!/[k!(n-k)!]$ counts success-position choices. $n!$ multiplies integers from 1 to $n$; $0!=1$. & \textbullet\ \hyperref[ex:intro-two]{Example 2}: $\binom32=3$ possible orders with two payments. \\[4pt]
$\GTenHLProbability{P(X=k)}$; $\HLCumulativeProbability{P(X\leq b)}$ & \GTenHLProbability{Event probability} & Probability of exactly $k$, or at most $b$, successes. Probabilities are unitless. & \textbullet\ \hyperref[ex:intro-two]{Example 2}: exactly two paid invoices has probability $\GTenHLProbability{0.375}$. \\[4pt]
$\sum_{j=a}^{b}$ & Summation & Add a term for each integer index $j$ from $a$ to $b$, including endpoints. & \textbullet\ \hyperref[ex:intro-two]{Example 2}: $\sum_{j=0}^{2}P(X=j)$ adds the probabilities of zero, one or two payments. \\[4pt]
$\GTenHLSecond{Y=5X}$ & \GTenHLSecond{Revenue} & A monetary consequence of a count, with 5 USD earned per success. & \textbullet\ \hyperref[ex:intro-one]{Example 1}: revenue in USD, distinct from the purchase count and its probability. \\
\bottomrule
\end{longtable}
\endgroup
\clearpage
\subsection{Practice procedure}
\begin{enumerate}
\item Define the random variable in words and give its units and support.
\item List parameters and check that the model assumptions match the situation.
\item Translate ordinary language into an event, paying attention to endpoints.
\item Write the symbolic formula first, then substitute without rounding intermediate values.
\item Check bounds, total probability or a complementary calculation; interpret the answer in context.
\end{enumerate}

\subsection{Bridge to the general model}
\[
X\sim\operatorname{Bin}(n,p),\qquad P(X=k)=\binom nk p^k(1-p)^{n-k},\quad k=0,\ldots,n.
\]
The powers describe one sequence: $k$ successes and $n-k$ failures. The combination counts its possible arrangements. For at most $b$, add terms from 0 to $b$. For at least $a$, use $1-P(X\leq a-1)$; the complement ends at $a-1$, not $a$. For between two values, state whether both endpoints are included.
\clearpage
\section{General Worked Problems --- Exact Counts}\label{sec:t2-c4-foundational}
\SectionProblemsTableOfContents{
\SectionProblemLink{prob:exactly-two}{Introductory general problem: exactly two successes}\\
\SectionProblemLink{prob:loans}{Loan defaults}\\
\SectionProblemLink{prob:invoices}{On-time invoices}
}

\Needspace{14\baselineskip}
\subsection{Introductory general problem: exactly two successes}\label{prob:exactly-two}
Three independent binary trials each have success probability $1/2$.

\QuestionsTasks[sec:t2-c4-foundational]{{4}/{sol:exactly-two}/{Calculate $P(X=2)$, defining success as the trial gives the designated S outcome.}}

\subsubsection*{Complete worked solution}\label{sol:exactly-two}
\begin{center}
\begin{tabular}{llp{0.48\linewidth}}
\toprule
Symbol & Value & Meaning in this problem \\ \midrule
$n$ & $3$ trials & Fixed number of trials \\
$p$ & $0.5$ & Success probability per trial \\
$k$ & $2$ successes & Requested success count \\
\bottomrule
\end{tabular}
\end{center}

\par\noindent\begin{minipage}{\linewidth}
\textbf{Step 1: Define the variable.} $X$ counts successes among 3 trials. A success means the trial gives the designated S outcome. Its support is $\{0,1,\ldots,3\}$ successes; money is a separate consequence.
\end{minipage}\par\medskip
\par\noindent\begin{minipage}{\linewidth}
\textbf{Step 2: State parameters and check conditions.} $\HLObservationCount{n}=\HLObservationCount{3}$ is fixed, $\HLSuccessProbability{p}=\HLSuccessProbability{0.5}$ and $\HLComplementProbability{1-p}=\HLComplementProbability{0.5}$. Under the stated independence, binary-outcome and constant-probability assumptions, $X\sim\operatorname{Bin}(3,0.5)$.
\end{minipage}\par\medskip
\par\noindent\begin{minipage}{\linewidth}
\textbf{Step 3: Translate and calculate.} The requested event is $ X=2$. First use
\[P(X=k)=\binom nk p^k(1-p)^{n-k}.\]
Then substitute:
\[P(X=2)=\binom{\HLObservationCount{3}}{\GTenHLValue{2}}(\HLSuccessProbability{0.5})^{2}(\HLComplementProbability{0.5})^{1}\approx \GTenHLProbability{0.3750}.\]
\end{minipage}\par\medskip
\par\noindent\begin{minipage}{\linewidth}
\textbf{Step 4: Check and interpret.} The answer is between 0 and 1; the full binomial PMF totals $(p+(1-p))^n=1$. Under this hypothetical model, the probability of the requested success count is $\GTenHLProbability{0.3750}$.
\end{minipage}\par\medskip
\secInicio[prob:exactly-two][sec:t2-c4-foundational]
\Needspace{14\baselineskip}
\Needspace{19\baselineskip}

\Needspace{14\baselineskip}
\clearpage
\subsection{Loan defaults}\label{prob:loans}
Five independent loans each default within one year with probability 0.1. Default is success mathematically.

\QuestionsTasks[sec:t2-c4-foundational]{{4}/{sol:loans}/{Calculate $P(X=1)$, defining success as a loan defaults within the year.}}

\subsubsection*{Complete worked solution}\label{sol:loans}
\par\noindent\begin{minipage}{\linewidth}
\textbf{Step 1: Define the variable.} $X$ counts successes among 5 trials. A success means a loan defaults within the year. Its support is $\{0,1,\ldots,5\}$ successes; money is a separate consequence.
\end{minipage}\par\medskip
\par\noindent\begin{minipage}{\linewidth}
\textbf{Step 2: State parameters and check conditions.} $\HLObservationCount{n}=\HLObservationCount{5}$ is fixed, $\HLSuccessProbability{p}=\HLSuccessProbability{0.1}$ and $\HLComplementProbability{1-p}=\HLComplementProbability{0.9}$. Under the stated independence, binary-outcome and constant-probability assumptions, $X\sim\operatorname{Bin}(5,0.1)$.
\end{minipage}\par\medskip
\par\noindent\begin{minipage}{\linewidth}
\textbf{Step 3: Translate and calculate.} The requested event is $ X=1$. First use
\[P(X=k)=\binom nk p^k(1-p)^{n-k}.\]
Then substitute:
\[P(X=1)=\binom{\HLObservationCount{5}}{\GTenHLValue{1}}(\HLSuccessProbability{0.1})^{1}(\HLComplementProbability{0.9})^{4}\approx \GTenHLProbability{0.3281}.\]
\end{minipage}\par\medskip
\par\noindent\begin{minipage}{\linewidth}
\textbf{Step 4: Check and interpret.} The answer is between 0 and 1; the full binomial PMF totals $(p+(1-p))^n=1$. Under this hypothetical model, the probability of the requested success count is $\GTenHLProbability{0.3281}$.
\end{minipage}\par\medskip

\secInicio[prob:loans][sec:t2-c4-foundational]
\Needspace{14\baselineskip}
\clearpage
\subsection{On-time invoices}\label{prob:invoices}
Four invoices are paid independently, each with constant probability 0.8 of being on time.

\QuestionsTasks[sec:t2-c4-foundational]{{4}/{sol:invoices}/{Calculate $P(X=3)$, defining success as an invoice is paid on time.}}

\subsubsection*{Complete worked solution}\label{sol:invoices}
\par\noindent\begin{minipage}{\linewidth}
\textbf{Step 1: Define the variable.} $X$ counts successes among 4 trials. A success means an invoice is paid on time. Its support is $\{0,1,\ldots,4\}$ successes; money is a separate consequence.
\end{minipage}\par\medskip
\par\noindent\begin{minipage}{\linewidth}
\textbf{Step 2: State parameters and check conditions.} $\HLObservationCount{n}=\HLObservationCount{4}$ is fixed, $\HLSuccessProbability{p}=\HLSuccessProbability{0.8}$ and $\HLComplementProbability{1-p}=\HLComplementProbability{0.2}$. Under the stated independence, binary-outcome and constant-probability assumptions, $X\sim\operatorname{Bin}(4,0.8)$.
\end{minipage}\par\medskip
\par\noindent\begin{minipage}{\linewidth}
\textbf{Step 3: Translate and calculate.} The requested event is $ X=3$. First use
\[P(X=k)=\binom nk p^k(1-p)^{n-k}.\]
Then substitute:
\[P(X=3)=\binom{\HLObservationCount{4}}{\GTenHLValue{3}}(\HLSuccessProbability{0.8})^{3}(\HLComplementProbability{0.2})^{1}\approx \GTenHLProbability{0.4096}.\]
\end{minipage}\par\medskip
\par\noindent\begin{minipage}{\linewidth}
\textbf{Step 4: Check and interpret.} The answer is between 0 and 1; the full binomial PMF totals $(p+(1-p))^n=1$. Under this hypothetical model, the probability of the requested success count is $\GTenHLProbability{0.4096}$.
\end{minipage}\par\medskip

\secInicio[prob:invoices][sec:t2-c4-foundational]
\clearpage
\section{General Worked Problems --- Cumulative Events}\label{sec:t2-c4-intermediate}
\SectionProblemsTableOfContents{
\SectionProblemLink{prob:interval}{A central interval}\\
\SectionProblemLink{prob:claims}{Insurance claims at most one}\\
\SectionProblemLink{prob:sales}{At least one business sale}
}

\subsection{A central interval}\label{prob:interval}
Six independent binary trials have constant success probability 0.4.

\QuestionsTasks[sec:t2-c4-intermediate]{{4}/{sol:interval}/{Calculate $P(2\leq X\leq4)$, defining success as the designated event occurs.}}

\subsubsection*{Complete worked solution}\label{sol:interval}
\par\noindent\begin{minipage}{\linewidth}
\textbf{Step 1: Define the variable.} $X$ counts successes among 6 trials. A success means the designated event occurs. Its support is $\{0,1,\ldots,6\}$ successes; money is a separate consequence.
\end{minipage}\par\medskip
\par\noindent\begin{minipage}{\linewidth}
\textbf{Step 2: State parameters and check conditions.} $\HLObservationCount{n}=\HLObservationCount{6}$ is fixed, $\HLSuccessProbability{p}=\HLSuccessProbability{0.4}$ and $\HLComplementProbability{1-p}=\HLComplementProbability{0.6}$. Under the stated independence, binary-outcome and constant-probability assumptions, $X\sim\operatorname{Bin}(6,0.4)$.
\end{minipage}\par\medskip
\par\noindent\begin{minipage}{\linewidth}
\textbf{Step 3: Translate and calculate.} The requested event is $ 2\leq X\leq4$. First use
\[P(a\leq X\leq b)=\sum_{j=a}^{b}\binom nj p^j(1-p)^{n-j}.\]
Then substitute:
\[P(2\leq X\leq4)=\sum_{j=2}^{4}\binom{\HLObservationCount{6}}j(\HLSuccessProbability{0.4})^j(\HLComplementProbability{0.6})^{6-j}\approx \GTenHLProbability{0.7258}.\]
\end{minipage}\par\medskip
\par\noindent\begin{minipage}{\linewidth}
\textbf{Step 4: Check and interpret.} The answer is between 0 and 1; the full binomial PMF totals $(p+(1-p))^n=1$. Under this hypothetical model, the probability of the requested success count is $\GTenHLProbability{0.7258}$.
\end{minipage}\par\medskip
\secInicio[prob:interval][sec:t2-c4-intermediate]
\Needspace{14\baselineskip}
\clearpage
\subsection{Insurance claims at most one}\label{prob:claims}
Six independent policies each have probability 0.05 of a claim this month.

\QuestionsTasks[sec:t2-c4-intermediate]{{4}/{sol:claims}/{Calculate $P(X\leq1)$, defining success as a policy has at least one claim this month.}}

\subsubsection*{Complete worked solution}\label{sol:claims}
\par\noindent\begin{minipage}{\linewidth}
\textbf{Step 1: Define the variable.} $X$ counts successes among 6 trials. A success means a policy has at least one claim this month. Its support is $\{0,1,\ldots,6\}$ successes; money is a separate consequence.
\end{minipage}\par\medskip
\par\noindent\begin{minipage}{\linewidth}
\textbf{Step 2: State parameters and check conditions.} $\HLObservationCount{n}=\HLObservationCount{6}$ is fixed, $\HLSuccessProbability{p}=\HLSuccessProbability{0.05}$ and $\HLComplementProbability{1-p}=\HLComplementProbability{0.95}$. Under the stated independence, binary-outcome and constant-probability assumptions, $X\sim\operatorname{Bin}(6,0.05)$.
\end{minipage}\par\medskip
\par\noindent\begin{minipage}{\linewidth}
\textbf{Step 3: Translate and calculate.} The requested event is $ X\leq1$. First use
\[P(a\leq X\leq b)=\sum_{j=a}^{b}\binom nj p^j(1-p)^{n-j}.\]
Then substitute:
\[P(X\leq1)=\sum_{j=0}^{1}\binom{\HLObservationCount{6}}j(\HLSuccessProbability{0.05})^j(\HLComplementProbability{0.95})^{6-j}\approx \GTenHLProbability{0.9672}.\]
\end{minipage}\par\medskip
\par\noindent\begin{minipage}{\linewidth}
\textbf{Step 4: Check and interpret.} The answer is between 0 and 1; the full binomial PMF totals $(p+(1-p))^n=1$. Under this hypothetical model, the probability of the requested success count is $\GTenHLProbability{0.9672}$.
\end{minipage}\par\medskip

\secInicio[prob:claims][sec:t2-c4-intermediate]
\Needspace{14\baselineskip}
\clearpage
\subsection{At least one business sale}\label{prob:sales}
Eight independent enquiries each have constant sale probability 0.25.

\QuestionsTasks[sec:t2-c4-intermediate]{{4}/{sol:sales}/{Calculate $P(X\geq1)$, defining success as an enquiry becomes a sale.}}

\subsubsection*{Complete worked solution}\label{sol:sales}
\par\noindent\begin{minipage}{\linewidth}
\textbf{Step 1: Define the variable.} $X$ counts successes among 8 trials. A success means an enquiry becomes a sale. Its support is $\{0,1,\ldots,8\}$ successes; money is a separate consequence.
\end{minipage}\par\medskip
\par\noindent\begin{minipage}{\linewidth}
\textbf{Step 2: State parameters and check conditions.} $\HLObservationCount{n}=\HLObservationCount{8}$ is fixed, $\HLSuccessProbability{p}=\HLSuccessProbability{0.25}$ and $\HLComplementProbability{1-p}=\HLComplementProbability{0.75}$. Under the stated independence, binary-outcome and constant-probability assumptions, $X\sim\operatorname{Bin}(8,0.25)$.
\end{minipage}\par\medskip
\par\noindent\begin{minipage}{\linewidth}
\textbf{Step 3: Translate and calculate.} The requested event is $ X\geq1$. First use
\[P(X\geq k)=1-\sum_{j=0}^{k-1}\binom nj p^j(1-p)^{n-j}.\]
Then substitute:
\[P(X\geq1)=1-\sum_{j=0}^{0}\binom{\HLObservationCount{8}}j(\HLSuccessProbability{0.25})^j(\HLComplementProbability{0.75})^{8-j}\approx \GTenHLProbability{0.8999}.\]
\end{minipage}\par\medskip
\par\noindent\begin{minipage}{\linewidth}
\textbf{Step 4: Check and interpret.} The answer is between 0 and 1; the full binomial PMF totals $(p+(1-p))^n=1$. Under this hypothetical model, the probability of the requested success count is $\GTenHLProbability{0.8999}$. For at least one, the complementary event is no successes.
\end{minipage}\par\medskip

\secInicio[prob:sales][sec:t2-c4-intermediate]
\clearpage
\section{Advanced Problems}\label{sec:t2-c4-advanced}
\SectionProblemsTableOfContents{
\SectionProblemLink{prob:fees}{Payment fees in a band}\\
\SectionProblemLink{prob:defects}{Supplier defects and a failed assumption}
}

\Needspace{14\baselineskip}
\subsection{Payment fees in a band}\label{prob:fees}
A firm processes ten independent payments, each late with probability 0.2. A late payment costs 5 USD, so total late fees are $Y=5X$ USD.

\QuestionsTasks[sec:t2-c4-advanced]{{4}/{sol:fees}/{Calculate $P(2\leq X\leq4)$, defining success as a payment is late. Explain why this is also the event $10\leq Y\leq20$ USD. Then change the fee to 10 USD per late payment: find the probability of the same monetary band and explain what changes in your model or event.}}

\subsubsection*{Complete worked solution}\label{sol:fees}
\par\noindent\begin{minipage}{\linewidth}
\textbf{Step 1: Define the variable.} $X$ counts successes among 10 trials. A success means a payment is late. Its support is $\{0,1,\ldots,10\}$ successes; money is a separate consequence.
\end{minipage}\par\medskip
\par\noindent\begin{minipage}{\linewidth}
\textbf{Step 2: State parameters and check conditions.} $\HLObservationCount{n}=\HLObservationCount{10}$ is fixed, $\HLSuccessProbability{p}=\HLSuccessProbability{0.2}$ and $\HLComplementProbability{1-p}=\HLComplementProbability{0.8}$. Under the stated independence, binary-outcome and constant-probability assumptions, $X\sim\operatorname{Bin}(10,0.2)$.
\end{minipage}\par\medskip
\par\noindent\begin{minipage}{\linewidth}
\textbf{Step 3: Translate and calculate.} The requested event is $ 2\leq X\leq4$. First use
\[P(a\leq X\leq b)=\sum_{j=a}^{b}\binom nj p^j(1-p)^{n-j}.\]
Then substitute:
\[P(2\leq X\leq4)=\sum_{j=2}^{4}\binom{\HLObservationCount{10}}j(\HLSuccessProbability{0.2})^j(\HLComplementProbability{0.8})^{10-j}\approx \GTenHLProbability{0.5914}.\]
\end{minipage}\par\medskip
\par\noindent\begin{minipage}{\linewidth}
\textbf{Step 4: Check and interpret.} The answer is between 0 and 1; the full binomial PMF totals $(p+(1-p))^n=1$. Under this hypothetical model, the probability of the requested success count is $\GTenHLProbability{0.5914}$. The fee transformation gives $10\leq5X\leq20$ USD, hence $2\leq X\leq4$. Both monetary endpoints are included.

\textbf{Follow-up: a changed fee.} If each late payment costs 10 USD instead, the same 10--20 USD band means $1\leq X\leq2$. The count model remains $X\sim\operatorname{Bin}(10,0.2)$, but the event changes:
\[P(1\leq X\leq2)=\binom{10}{1}(0.2)(0.8)^9+\binom{10}{2}(0.2)^2(0.8)^8\approx \GTenHLProbability{0.5704}.\]
The probability falls from $\GTenHLProbability{0.5914}$ to $\GTenHLProbability{0.5704}$ for this particular band; doubling the fee does not double the probability.
\end{minipage}\par\medskip

\secInicio[prob:fees][sec:t2-c4-advanced]
\Needspace{14\baselineskip}
\clearpage
\subsection{Supplier defects and a failed assumption}\label{prob:defects}
In version A, seven independently produced items each have defect probability 0.15. In version B, seven items are drawn without replacement from a small lot containing both defective and sound items.

\QuestionsTasks[sec:t2-c4-advanced]{{4}/{sol:defects}/{Calculate $P(X\geq2)$, defining success as an item is defective. Calculate for version A; explain why version B does not justify the same binomial model.}}

\subsubsection*{Complete worked solution}\label{sol:defects}
\par\noindent\begin{minipage}{\linewidth}
\textbf{Step 1: Define the variable.} $X$ counts successes among 7 trials. A success means an item is defective. Its support is $\{0,1,\ldots,7\}$ successes; money is a separate consequence.
\end{minipage}\par\medskip
\par\noindent\begin{minipage}{\linewidth}
\textbf{Step 2: State parameters and check conditions.} $\HLObservationCount{n}=\HLObservationCount{7}$ is fixed, $\HLSuccessProbability{p}=\HLSuccessProbability{0.15}$ and $\HLComplementProbability{1-p}=\HLComplementProbability{0.85}$. Under the stated independence, binary-outcome and constant-probability assumptions, $X\sim\operatorname{Bin}(7,0.15)$.
\end{minipage}\par\medskip
\par\noindent\begin{minipage}{\linewidth}
\textbf{Step 3: Translate and calculate.} The requested event is $ X\geq2$. First use
\[P(X\geq k)=1-\sum_{j=0}^{k-1}\binom nj p^j(1-p)^{n-j}.\]
Then substitute:
\[P(X\geq2)=1-\sum_{j=0}^{1}\binom{\HLObservationCount{7}}j(\HLSuccessProbability{0.15})^j(\HLComplementProbability{0.85})^{7-j}\approx \GTenHLProbability{0.2834}.\]
\end{minipage}\par\medskip
\par\noindent\begin{minipage}{\linewidth}
\textbf{Step 4: Check and interpret.} The answer is between 0 and 1; the full binomial PMF totals $(p+(1-p))^n=1$. Under this hypothetical model, the probability of the requested success count is $\GTenHLProbability{0.2834}$. In the second version, sampling without replacement from a small lot changes the defect probability after each draw. The binomial assumptions fail; a different sampling model or the full lot composition is needed.
\end{minipage}\par\medskip

\secInicio[prob:defects][sec:t2-c4-advanced]
\subsection*{Key ideas and common errors}
Check fixed $n$, binary outcomes, independence and constant $p$ before calculating. The combination counts arrangements; do not omit it. At least $k$ complements at most $k-1$. Count successes separately from monetary costs. Sampling without replacement and common shocks can invalidate the model.
\secInicio
\end{document}
